\scanpage{75}
- \mspage{210} -
\margin{n$^\circ$ 391}
\[
\begin{aligned}
\alpha_1&=\tfrac12(\varepsilon_1+\varepsilon_8)-\tfrac12(\varepsilon_2+\varepsilon_3+\varepsilon_4+\varepsilon_5+\varepsilon_6+\varepsilon_7),\quad
\alpha_2=\varepsilon_1+\varepsilon_2,\\
\alpha_3&=\varepsilon_2-\varepsilon_1,\quad
\alpha_4=\varepsilon_3-\varepsilon_2,\quad
\alpha_5=\varepsilon_4-\varepsilon_3,\quad
\alpha_6=\varepsilon_5-\varepsilon_4,\quad
\alpha_7=\varepsilon_6-\varepsilon_5,\\
\alpha_8&=\varepsilon_7-\varepsilon_6.
\end{aligned}
\]
Il est clair, en vertu de (I), que ce sont des racines de $R$, et qu'elles sont $>0$ pour l'ordre considéré. D'autre part, pour cet ordre, les racines positives sont les $\varepsilon_j\pm\varepsilon_i$ pour $i<j$ et les racines $\tfrac12(\varepsilon_8+\sum_{i=1}^7(-1)^{\unsure{\vartheta|\delta}(i)}\varepsilon_i)$ avec $\sum_{i=1}^7\unsure{\vartheta|\delta}(i)$ pair. Montrons que ces racines sont combinaisons linéaires des $\alpha_i$ à coefficients tous $>0$. Pour les racines $\varepsilon_j\pm\varepsilon_i$ avec $1\leq i<j\leq7$, cela résulte du calcul fait pour le système de type $D_7$, (n$^\circ$ 8). Il en est de même de $\varepsilon_1+\varepsilon_8$ par l'expression de $\alpha_1$, puis, par récurrence sur $i$, de $\varepsilon_8+\varepsilon_i=(\varepsilon_8+\varepsilon_{i-1})+(\varepsilon_i-\varepsilon_{i-1})$ pour $2\leq i\leq7$. En considérant ensuite l'expression $2\alpha_1+\alpha_3$, on obtient $\varepsilon_8-\varepsilon_7$ comme combinaison des $\alpha_i$ à coefficients $\geq0$, puis par récurrence descendante sur $i$, $\varepsilon_8-\varepsilon_i=(\varepsilon_8-\varepsilon_{i+1})+(\varepsilon_{i+1}-\varepsilon_i)$ pour $1\leq i\leq6$. En ajoutant à $\alpha_1$ des éléments $\varepsilon_i+\varepsilon_j$ avec $2\leq i<j\leq7$, on obtient toutes les racines de la forme $\tfrac12(\varepsilon_1+\varepsilon_8)+\tfrac12(\sum_{i=2}^7(-1)^{\unsure{\vartheta|\delta}(i)}\varepsilon_i)$ avec $\sum_{i=2}^7\unsure{\vartheta|\delta}(i)$ pair ; enfin, $\alpha_1+\alpha_2$ a la même forme que $\alpha_1$ sauf qu'on y a permuté $\varepsilon_1$ et $\varepsilon_2$, et en lui ajoutant des éléments $\varepsilon_i+\varepsilon_j$ avec $1\leq i<j\leq7$, $i$, $j$ distincts de $2$, on achève de prouver notre assertion. Les $\alpha_i$, $(1\leq i\leq8)$, forment donc un système de \uline{racines simples} $S$ pour $R$, (\S\ 6, n$^\circ$ 7, cor. 2 de la prop. 14). En outre, on a $(\alpha_i|\alpha_i)=2$ pour tout $i$, $(\alpha_i|\alpha_{i+1})=(\alpha_2|\alpha_4)=(\alpha_1|\alpha_3)=-1$, $(3\leq i\leq8)$, et $(\alpha_i|\alpha_j)=0$ pour tous les autres couples d'indices ; le schéma de Dynkin de $R$ est donc de type $E_8$, et $R$ est irréductible.
\mspage{74}
